\documentclass[12pt,a4paper]{article}

% ── Packages ──────────────────────────────────────────────────────────────────
\usepackage[a4paper, margin=2.5cm]{geometry}
\usepackage{fontenc}
\usepackage{inputenc}
\usepackage{microtype}
\usepackage{lmodern}
\usepackage{hyperref}
\usepackage{graphicx}
\usepackage{float}
\usepackage{booktabs}
\usepackage{tabularx}
\usepackage{multirow}
\usepackage{array}
\usepackage{xcolor}
\usepackage{listings}
\usepackage{enumitem}
\usepackage{titlesec}
\usepackage{fancyhdr}
\usepackage{caption}
\usepackage{subcaption}
\usepackage{parskip}
\usepackage{amsmath}

% ── Colors ────────────────────────────────────────────────────────────────────
\definecolor{accent}{HTML}{0f3460}
\definecolor{lightgray}{HTML}{f4f6f9}
\definecolor{codegray}{HTML}{f0f0f0}
\definecolor{darkgray}{HTML}{444444}

% ── Hyperref setup ────────────────────────────────────────────────────────────
\hypersetup{
  colorlinks=true,
  linkcolor=accent,
  urlcolor=accent,
  citecolor=accent,
  pdftitle={Na'vi Translator -- MLOps Project Report},
  pdfauthor={DA6360 Student},
}

% ── Section title formatting ──────────────────────────────────────────────────
\titleformat{\section}
  {\large\bfseries\color{accent}}
  {\thesection.}{0.6em}{}[\titlerule]

\titleformat{\subsection}
  {\normalsize\bfseries\color{darkgray}}
  {\thesubsection}{0.5em}{}

% ── Header / Footer ───────────────────────────────────────────────────────────
\pagestyle{fancy}
\fancyhf{}
\rhead{\small Na'vi Translator --- DA6360 MLOps}
\lhead{\small \leftmark}
\rfoot{\small Page \thepage}
\renewcommand{\headrulewidth}{0.4pt}

% ── Code listing style ────────────────────────────────────────────────────────
\lstset{
  backgroundcolor=\color{codegray},
  basicstyle=\ttfamily\footnotesize,
  breaklines=true,
  frame=single,
  framesep=3pt,
  rulecolor=\color{lightgray},
  xleftmargin=8pt,
}

% ── Table helpers ─────────────────────────────────────────────────────────────
\newcolumntype{L}[1]{>{\raggedright\arraybackslash}p{#1}}
\newcolumntype{C}[1]{>{\centering\arraybackslash}p{#1}}

% ── Image placeholder ─────────────────────────────────────────────────────────
% Usage: \imgplaceholder{width}{height}{label}{caption}
\newcommand{\imgplaceholder}[4]{%
  \begin{figure}[H]
    \centering
    \fbox{\begin{minipage}[c][#2][c]{#1}
      \centering
      \textcolor{darkgray}{\textit{[Screenshot: #3]}}
    \end{minipage}}
    \caption{#4}
    \label{fig:#3}
  \end{figure}%
}

% ── Real image command ────────────────────────────────────────────────────────
% Usage: \img{width}{path}{label}{caption}
\newcommand{\img}[4]{%
  \begin{figure}[H]
    \centering
    \includegraphics[width=#1]{#2}
    \caption{#4}
    \label{fig:#3}
  \end{figure}%
}

% ──────────────────────────────────────────────────────────────────────────────
\begin{document}

% ── Title Page ────────────────────────────────────────────────────────────────
\begin{titlepage}
  \centering
  \vspace*{3cm}
  {\Huge\bfseries\color{accent} Na'vi Language Translator\par}
  \vspace{0.5cm}
  {\Large End-to-End MLOps Application\par}
  \vspace{1.5cm}
  \rule{0.6\textwidth}{0.4pt}
  \vspace{0.8cm}

  \begin{tabular}{ll}
    \textbf{Course}      & DA6360 --- MLOps \\[4pt]
    \textbf{Repository}  & \url{https://github.com/usnaveen/navi-translator} \\[4pt]
    \textbf{Submitted}   & 28 April 2026 \\
  \end{tabular}

  \vspace{2cm}

  \begin{tabular}{@{}L{5cm}L{8cm}@{}}
    \toprule
    \textbf{Item} & \textbf{Location} \\
    \midrule
    GitHub Repository  & github.com/usnaveen/navi-translator \\
    Architecture (HLD) & docs/HLD.md \\
    API Spec (LLD)     & docs/LLD.md \\
    Test Plan          & docs/TEST\_PLAN.md \\
    Test Report        & docs/TEST\_REPORT.md \\
    User Manual        & docs/user\_manual.md \\
    Automated Tests    & 32 / 32 passing \\
    \bottomrule
  \end{tabular}

\end{titlepage}

\tableofcontents
\newpage

% ──────────────────────────────────────────────────────────────────────────────
\section{Executive Summary}

Na'vi Translator is a production-grade MLOps application that translates Na'vi --- a
low-resource constructed language from the \textit{Avatar} film universe --- into English
from both audio and text input. The system implements the full MLOps lifecycle:
data ingestion $\rightarrow$ versioning $\rightarrow$ training $\rightarrow$ experiment
tracking $\rightarrow$ packaging $\rightarrow$ deployment $\rightarrow$ monitoring
$\rightarrow$ automated retraining.

The core engineering challenge is \textbf{low-resource adaptation under a strict
no-cloud constraint}. Na'vi has approximately 2,600 documented word pairs and limited
public audio data. All training runs on local Apple Silicon (MPS) hardware using
PEFT/LoRA to make Whisper fine-tuning feasible within the hardware constraint.

\subsection*{Key Achievements}
\begin{itemize}[leftmargin=1.5em]
  \item Full MLOps pipeline: DVC data versioning, Airflow orchestration, MLflow experiment
        tracking, Prometheus/Grafana monitoring
  \item Two fine-tuned models registered in MLflow Model Registry: Whisper-tiny + LoRA
        (ASR) and MarianMT (NMT)
  \item FastAPI backend with 6 REST endpoints; nginx frontend; Docker Compose with
        5 isolated services
  \item 32\,/\,32 automated tests passing (unit + integration); GitHub Actions CI/CD
  \item Reykunyu dictionary fallback for outputs below the confidence threshold
\end{itemize}

% ──────────────────────────────────────────────────────────────────────────────
\section{Problem Statement and Success Metrics}

Na'vi has an active hobbyist learning community but no unified translation tool.
Static dictionary lookups exist but cannot handle audio, multi-word phrases, or
community vocabulary contributions. The goal is a single system covering speech,
text, and community submissions with full MLOps observability.

\begin{table}[H]
  \centering
  \caption{Project success metrics}
  \begin{tabular}{@{}L{2cm}L{6cm}L{2.5cm}L{3.5cm}@{}}
    \toprule
    \textbf{Type} & \textbf{Metric} & \textbf{Target} & \textbf{Result} \\
    \midrule
    ML       & Whisper ASR Word Error Rate  & $\leq 0.35$ & 0.42 (low-resource) \\
    ML       & MarianMT BLEU score          & $\geq 0.25$ & Logged in MLflow   \\
    Business & Text translation latency     & $<$ 200\,ms & $\sim$85\,ms        \\
    Business & Audio translation latency    & $<$ 5\,s    & $\sim$1.2\,s        \\
    Business & OOV drift alert threshold    & $<$ 15\%    & Prometheus alert configured \\
    \bottomrule
  \end{tabular}
\end{table}

% ──────────────────────────────────────────────────────────────────────────────
\section{System Architecture}

The system is split into four layers --- Data, Training, Serving, and Monitoring ---
each independently containerised and connected through defined interfaces.

\begin{table}[H]
  \centering
  \caption{System layers and components}
  \begin{tabular}{@{}L{2.2cm}L{7cm}L{5cm}@{}}
    \toprule
    \textbf{Layer} & \textbf{Components} & \textbf{Tools} \\
    \midrule
    Data      & Ingestion, versioning, preprocessing, splitting, baseline stats
              & DVC, Reykunyu API, gTTS, Airflow \\
    Training  & Whisper-tiny + LoRA (ASR), MarianMT (NMT), evaluation, auto-promotion
              & MLflow, PEFT, HuggingFace Transformers \\
    Serving   & REST API, model inference, fallback dictionary, metrics export
              & FastAPI, nginx, Prometheus client \\
    Monitoring& Metric scraping, dashboards, drift detection, retrain alerts
              & Prometheus, Grafana, EvidentlyAI \\
    \bottomrule
  \end{tabular}
\end{table}

\subsection{Request Flow --- Audio Translation}

\begin{lstlisting}
Browser MediaRecorder
  --> POST /api/translate/audio
  --> Resample to 16 kHz mono (librosa)
  --> Whisper-tiny encoder + LoRA-adapted decoder
  --> Na'vi text
  --> MarianMT NMT
  --> English translation
  --> if confidence < 0.4: Reykunyu word-by-word fallback
  --> JSON response {navi_text, english, confidence, latency_ms}
\end{lstlisting}

\subsection{Docker Compose Network}

\begin{lstlisting}
navi-net (bridge)
  +-- backend    :8000  FastAPI inference server
  +-- frontend   :3000  nginx static UI + reverse proxy
  +-- mlflow     :5000  MLflow tracking server
  +-- prometheus :9090  metric scraper
  +-- grafana    :3001  near-real-time dashboard
\end{lstlisting}

% ──────────────────────────────────────────────────────────────────────────────
\section{MLOps Implementation}

\subsection{Data Engineering}

The data pipeline is defined as a DVC DAG in \texttt{dvc.yaml} with five stages.
Each stage declares source-code dependencies, parameter dependencies (from
\texttt{params.yaml}), and output artifacts so DVC detects changes and re-runs
only what is necessary.

\begin{table}[H]
  \centering
  \caption{DVC pipeline stages}
  \begin{tabular}{@{}L{3cm}L{4.5cm}L{4cm}L{2.5cm}@{}}
    \toprule
    \textbf{Stage} & \textbf{Input} & \textbf{Output} & \textbf{Tool} \\
    \midrule
    ingest           & Reykunyu API + TTS     & words.json + audio  & requests, gTTS \\
    preprocess\_text & words.json             & processed/text      & pandas         \\
    preprocess\_audio& raw audio              & processed/audio     & librosa        \\
    split            & text + audio           & train/val/test      & sklearn        \\
    baseline\_stats  & train split            & baselines.json      & numpy          \\
    \bottomrule
  \end{tabular}
\end{table}

\begin{itemize}[leftmargin=1.5em]
  \item 3,956 training examples, 494 validation examples; full pipeline completes in under 5 minutes
  \item Baseline statistics (mean, variance, OOV rate) stored in \texttt{baselines.json} for drift comparison
  \item Apache Airflow DAG orchestrates the same pipeline in production with task-level retry
\end{itemize}

% ── IMAGE 1: DVC DAG ──────────────────────────────────────────────────────────
\img{0.85\textwidth}{screenshots/dvc-dag.png}{dvc-dag}
  {DVC pipeline DAG showing ingest, preprocessing, split, and baseline stages}

\subsection{Source Control and Continuous Integration}

\begin{itemize}[leftmargin=1.5em]
  \item \textbf{Git} --- source code, configuration, and documentation
  \item \textbf{DVC} --- large data artifacts and model weights (pointer files in Git,
        content in \texttt{.dvc/cache})
  \item \textbf{GitHub Actions} --- CI on every push: \texttt{ruff} lint
        $\rightarrow$ unit tests $\rightarrow$ integration tests $\rightarrow$ docker build
  \item Every MLflow run is tagged with its Git commit SHA and DVC lock hash for full
        reproducibility: a run is reproducible from \texttt{(git\_sha, dvc\_lock\_sha, mlflow\_run\_id)}
\end{itemize}

\subsection{Experiment Tracking (MLflow)}

Every training run is tracked in MLflow with full parameter, metric, and artifact logging.

\begin{table}[H]
  \centering
  \caption{MLflow tracked items}
  \begin{tabular}{@{}L{2.5cm}L{11.5cm}@{}}
    \toprule
    \textbf{Item} & \textbf{Details} \\
    \midrule
    Parameters & base\_model, lora\_r, lora\_alpha, lora\_target\_modules, learning\_rate,
                 batch\_size, num\_epochs, warmup\_steps \\
    Metrics    & train/loss every 25 steps, eval/WER, eval/BLEU, final\_wer, final\_bleu \\
    Artifacts  & LoRA adapter weights, tokenizer, processor, training\_args.bin \\
    Tags       & git\_sha, dvc\_lock\_sha --- full reproducibility guarantee \\
    Registry   & navi-whisper (v1) and navi-marian (v1) in Production stage;
                 auto-promote if WER/BLEU improves by $\geq 0.01$ \\
    \bottomrule
  \end{tabular}
\end{table}

% ── IMAGE 3: MLflow loss curve ────────────────────────────────────────────────
\img{\textwidth}{screenshots/mlflow-loss-curve.png}{mlflow-loss-curve}
  {Training loss curve from the final whisper-tiny run (loss: 5.12 $\rightarrow$ 0.86)}

% ── IMAGE 4: MLflow model registry ───────────────────────────────────────────
\img{0.8\textwidth}{screenshots/mlflow-registry.png}{mlflow-registry}
  {MLflow Model Registry showing navi-whisper and navi-marian in Production stage}

\subsection{Exporter Instrumentation and Visualization}

The FastAPI backend exposes \texttt{/metrics} in Prometheus text format.
Prometheus scrapes it every 10 seconds; Grafana dashboards visualise all metrics
in near-real-time.

\begin{table}[H]
  \centering
  \caption{Prometheus metrics}
  \begin{tabular}{@{}L{5.5cm}L{2cm}L{6.5cm}@{}}
    \toprule
    \textbf{Metric} & \textbf{Type} & \textbf{Purpose} \\
    \midrule
    navi\_translation\_requests\_total  & Counter   & Request volume by input type and status \\
    navi\_translation\_latency\_ms      & Histogram & End-to-end latency distribution \\
    navi\_whisper\_wer\_live            & Gauge     & Live ASR word error rate \\
    navi\_marian\_bleu\_live            & Gauge     & Live NMT BLEU score \\
    navi\_oov\_rate                     & Gauge     & Out-of-vocabulary rate (primary drift signal) \\
    navi\_fallback\_rate                & Gauge     & Fraction of requests using dictionary fallback \\
    navi\_model\_version                & Info      & Currently loaded model version \\
    navi\_vocab\_submissions\_total     & Counter   & Community vocabulary submissions \\
    \bottomrule
  \end{tabular}
\end{table}

\textbf{Alert rules:}
\begin{itemize}[leftmargin=1.5em]
  \item \texttt{HighErrorRate} --- error rate $>$ 5\% over 5 minutes, severity critical
  \item \texttt{OOVDriftDetected} --- \texttt{navi\_oov\_rate} $>$ 0.15, triggers Airflow retrain DAG
  \item \texttt{HighLatency} --- p95 latency $>$ 5,000\,ms over 5 minutes, severity warning
\end{itemize}

% ── IMAGE 5: Grafana dashboard preview ───────────────────────────────────────
\img{0.42\textwidth}{screenshots/grafana-dashboard.png}{grafana-dashboard}
  {Grafana-style dashboard preview derived from the checked-in panel configuration for request volume, latency, OOV rate, and model version monitoring}

% ── IMAGE 6: Prometheus targets ───────────────────────────────────────────────
\img{0.8\textwidth}{screenshots/prometheus-targets.png}{prometheus-targets}
  {Prometheus scrape targets --- backend:8000/metrics showing UP status}

\subsection{Software Packaging}

\begin{itemize}[leftmargin=1.5em]
  \item \textbf{MLproject} --- defines \texttt{train\_whisper}, \texttt{train\_marian}, and
        \texttt{evaluate} entry points; \texttt{python\_env.yaml} pins the full dependency set
        for reproducible runs
  \item \textbf{FastAPI} --- REST inference API; Pydantic-enforced schemas; model loaded
        from MLflow registry at startup; \texttt{/health} and \texttt{/ready} endpoints for
        orchestration health checks
  \item \textbf{Docker Compose} --- 5 isolated services on a \texttt{navi-net} bridge network,
        each with its own Dockerfile and health check
  \item \textbf{MLflow Model Registry} --- models promoted to the Production alias; FastAPI loads
        the Production model at startup without code changes
\end{itemize}

% ──────────────────────────────────────────────────────────────────────────────
\section{Model Development and Training}

\subsection{Whisper-tiny + LoRA (ASR)}

Whisper-tiny (39M parameters) was selected over Whisper-small (244M) based on
published research \cite{springer2024} showing that smaller Whisper variants outperform
larger ones when fine-tuning data is limited. PEFT/LoRA adapters reduce the trainable
parameter count to 589,824 (1.54\% of total), making training feasible on Apple
Silicon MPS hardware within the no-cloud constraint.

\begin{table}[H]
  \centering
  \caption{Whisper-tiny + LoRA training configuration}
  \begin{tabular}{@{}L{3.5cm}L{4.5cm}L{6cm}@{}}
    \toprule
    \textbf{Parameter} & \textbf{Value} & \textbf{Rationale} \\
    \midrule
    Base model         & openai/whisper-tiny          & 39M params; best low-resource performance \cite{springer2024} \\
    LoRA rank          & 16                           & Bumped from 8 to compensate for smaller backbone \\
    LoRA alpha         & 32                           & Standard $\alpha = 2r$ ratio \\
    Target modules     & q\_proj, k\_proj, v\_proj, out\_proj & Full attention coverage \\
    Learning rate      & $1 \times 10^{-4}$           & Conservative; higher LR can hurt tiny variants \\
    Effective batch    & 4 (1 device $\times$ 4 accum)& Gradient accumulation avoids OOM \\
    Epochs             & 5                            & Full corpus sweep \\
    Trainable params   & 589,824 / 38,350,464 (1.54\%)& LoRA efficiency \\
    Training runtime   & 3373\,s ($\approx$56\,min)   & Apple M-series MPS; no cloud required \\
    Final train loss   & 0.38 (started 5.12)          & Good convergence over 4945 steps \\
    Language           & \texttt{mi} (M\={a}ori)      & Closest Whisper-supported language to Na'vi phonology \\
    Final WER          & 0.42                         & Target was $\leq 0.35$; limited by small audio corpus \\
    \bottomrule
  \end{tabular}
\end{table}

\subsection{MarianMT (Na'vi $\rightarrow$ English NMT)}

Helsinki-NLP/opus-mt-mul-en was selected as the base translation model --- a lightweight
seq2seq architecture purpose-built for machine translation that outperforms general-purpose
LLMs on domain-specific fine-tuning with small datasets.

\begin{table}[H]
  \centering
  \caption{MarianMT training configuration}
  \begin{tabular}{@{}L{4cm}L{10cm}@{}}
    \toprule
    \textbf{Parameter} & \textbf{Value} \\
    \midrule
    Base model    & Helsinki-NLP/opus-mt-mul-en \\
    Learning rate & $2 \times 10^{-5}$ \\
    Batch size    & 16 \\
    Epochs        & 10 \\
    Final train loss & $\approx 3.30$ \\
    Registry      & navi-marian v1, Production stage \\
    \bottomrule
  \end{tabular}
\end{table}

\subsection{Reykunyu Dictionary Fallback}

When neural translation confidence falls below 0.4, the system falls back to
word-by-word lookup in the Reykunyu Na'vi dictionary ($\sim$2,600 entries). This
keeps the application usable for common vocabulary even when the neural model
confidence is low --- a deliberate design choice for low-resource ASR.

% ──────────────────────────────────────────────────────────────────────────────
\section{Software Engineering}

\subsection{API Endpoints}

Full endpoint specifications including error codes are in \texttt{docs/LLD.md}.

\begin{table}[H]
  \centering
  \caption{REST API endpoint summary}
  \begin{tabular}{@{}C{1.2cm}L{3.5cm}L{4cm}L{5.3cm}@{}}
    \toprule
    \textbf{Method} & \textbf{Endpoint} & \textbf{Input} & \textbf{Output} \\
    \midrule
    POST & /translate/audio & multipart/form-data (WAV/MP3/OGG $\leq$ 30s) & navi\_text, english, confidence, latency\_ms \\
    POST & /translate/text  & JSON: \{text: 1--500 chars\} & english, confidence, word\_breakdown, latency\_ms \\
    POST & /vocab/submit    & JSON: \{navi\_word, english\_meaning, audio\_b64?\} & accepted, message \\
    GET  & /health          & ---  & status, model\_version, uptime\_s \\
    GET  & /ready           & ---  & ready, whisper\_loaded, marian\_loaded \\
    GET  & /metrics         & ---  & Prometheus text exposition \\
    \bottomrule
  \end{tabular}
\end{table}

\subsection{API Documentation and Health Checks}

The FastAPI service also exposes interactive OpenAPI documentation and lightweight
health endpoints for readiness checks during local deployment.

\img{0.42\textwidth}{screenshots/api-docs.png}{api-docs}
  {Swagger / OpenAPI documentation view for the Na'vi Translator backend}

\img{0.36\textwidth}{screenshots/api-health.png}{api-health}
  {Health endpoint response confirming service uptime and model version reporting}

\subsection{Implementation Quality}

\begin{itemize}[leftmargin=1.5em]
  \item \textbf{Code style:} PEP-8 enforced via \texttt{ruff}; linting runs on every CI push
  \item \textbf{Logging:} Module-level loggers throughout; training logs to stdout and timestamped log files
  \item \textbf{Exception handling:} FastAPI global handlers return structured JSON errors with request IDs
  \item \textbf{Type safety:} Pydantic v2 enforces all request/response schemas at runtime
  \item \textbf{Loose coupling:} Frontend (nginx + static JS) and backend communicate only via REST; API base URL is configurable
\end{itemize}

\subsection{Testing}

\begin{table}[H]
  \centering
  \caption{Automated test results}
  \begin{tabular}{@{}L{3.5cm}C{1.5cm}L{6.5cm}C{2.5cm}@{}}
    \toprule
    \textbf{Suite} & \textbf{Count} & \textbf{Coverage} & \textbf{Result} \\
    \midrule
    Unit tests        & 30 & Schemas, drift detection, Reykunyu parser, Prometheus helpers, build\_pairs & 30/30 PASSED \\
    Integration tests & 2  & FastAPI route wiring (fake engine, no model weights) & 2/2 PASSED \\
    \midrule
    \textbf{Total}    & \textbf{32} & & \textbf{32/32 PASSED} \\
    \bottomrule
  \end{tabular}
\end{table}

% ──────────────────────────────────────────────────────────────────────────────
\section{Web Application and Pipeline Visualization}

\subsection{Frontend UI}

A 4-tab single-page application served by nginx. Pure HTML/CSS/JavaScript ---
no framework dependencies --- for minimal footprint and easy containerisation.

\begin{itemize}[leftmargin=1.5em]
  \item \textbf{Translate Audio:} Browser \texttt{MediaRecorder} API captures microphone input or accepts file upload (WAV/MP3/OGG). Displays Na'vi transcription, English translation, and confidence score.
  \item \textbf{Translate Text:} Text input with Na'vi-to-English translation and word-level breakdown.
  \item \textbf{Submit Vocabulary:} Community contributions with optional audio upload; validated against Na'vi character set.
  \item \textbf{About:} System description and user manual.
\end{itemize}

% ── IMAGE 7: Frontend home ────────────────────────────────────────────────────
\img{\textwidth}{screenshots/frontend-home.png}{frontend-home}
  {Na'vi Translator web application --- 4-tab UI loaded at \texttt{http://localhost:3000}}

% ── IMAGE 8: Translation result ───────────────────────────────────────────────
\img{\textwidth}{screenshots/frontend-translation.png}{frontend-translation}
  {Text translation result for \textit{oel ngati kameie} showing word-level breakdown}

\subsection{MLOps Pipeline Visualization}

\begin{table}[H]
  \centering
  \caption{MLOps visualization tools accessible via Docker Compose}
  \begin{tabular}{@{}L{2.5cm}L{4cm}L{7.5cm}@{}}
    \toprule
    \textbf{Tool} & \textbf{URL} & \textbf{Purpose} \\
    \midrule
    MLflow UI  & \texttt{localhost:5000} & Experiment runs, loss curves, parameter comparison, model registry \\
    Grafana    & \texttt{localhost:3001} & Real-time inference metrics, OOV drift, latency histograms \\
    Prometheus & \texttt{localhost:9090} & Raw metric queries, alert state, scrape target health \\
    Airflow    & DAG CLI / UI           & Training DAG runs, task logs, retrain scheduling \\
    DVC CLI    & \texttt{dvc dag}       & Data pipeline DAG, stage lineage and caching \\
    \bottomrule
  \end{tabular}
\end{table}

% ──────────────────────────────────────────────────────────────────────────────
\section{Problems Faced and Mitigations}

\begin{table}[H]
  \centering
  \caption{Engineering challenges and resolutions}
  \begin{tabular}{@{}C{0.4cm}L{3.5cm}L{4.5cm}L{5.6cm}@{}}
    \toprule
    \textbf{\#} & \textbf{Problem} & \textbf{Root Cause} & \textbf{Mitigation} \\
    \midrule
    1 & Training crash wiped all progress
      & \texttt{save\_strategy=epoch} --- crash 1 step before epoch end
      & Switched to \texttt{save\_strategy=steps} with \texttt{save\_steps=500}; max 8 min lost per crash \\
    2 & Whisper-small training infeasible (4+ hours)
      & 244M params $\times$ $\sim$3 s/step on Apple Silicon MPS
      & Whisper-tiny (39M params, $\sim$6$\times$ faster); validated by Springer 2024 research for low-resource data \\
    3 & MPS thermal throttling during mid-training eval (240 s/step)
      & Mac throttles under sustained eval loops even without Docker contention
      & Removed mid-training eval (\texttt{eval\_strategy=no}); single eval runs once after all training \\
    4 & CI pipeline failing on every push
      & CI was installing torch/transformers/peft ($\sim$3\,GB) on every GitHub Actions run
      & Created \texttt{requirements-ci.txt} with only the 12 lightweight packages tests actually import \\
    5 & Limited Na'vi audio data ($\sim$4,000 samples)
      & Constructed language; no commercial corpus
      & Synthetic TTS augmentation + LoRA to avoid overfitting + Reykunyu fallback for low-confidence outputs \\
    6 & No-cloud constraint
      & Course guidelines explicitly prohibit cloud platforms
      & Whisper-tiny + LoRA makes training feasible locally; documented design choice aligned with guidelines \S{}II.C \\
    \bottomrule
  \end{tabular}
\end{table}

% ──────────────────────────────────────────────────────────────────────────────
\section{Technology Stack}

\begin{table}[H]
  \centering
  \caption{Full technology stack}
  \begin{tabular}{@{}L{3cm}L{4cm}L{2.5cm}L{4.5cm}@{}}
    \toprule
    \textbf{Category} & \textbf{Tool} & \textbf{Version} & \textbf{Role} \\
    \midrule
    ASR              & OpenAI Whisper-tiny + PEFT LoRA & whisper-tiny    & Speech-to-Na'vi-text \\
    NMT              & MarianMT                        & opus-mt-mul-en  & Na'vi-to-English translation \\
    API              & FastAPI + Uvicorn               & $\geq$0.109     & REST inference server \\
    Frontend         & nginx                           & alpine          & Static UI + reverse proxy \\
    Experiment Track & MLflow                          & $\geq$2.10      & Run tracking + model registry \\
    Data Version     & DVC                             & 3.x             & Data \& artifact versioning \\
    Orchestration    & Apache Airflow                  & 2.8+            & Training DAG scheduling \\
    Monitoring       & Prometheus + Grafana            & latest          & Metrics scraping + dashboards \\
    Containers       & Docker Compose                  & v2              & 5-service local deployment \\
    CI/CD            & GitHub Actions                  & ---             & Lint + test + docker build \\
    Testing          & pytest + pytest-cov             & $\geq$9.0       & Unit + integration tests \\
    Linting          & ruff                            & latest          & PEP-8 enforcement \\
    \bottomrule
  \end{tabular}
\end{table}

% ──────────────────────────────────────────────────────────────────────────────
\section{Conclusion}

The Na'vi Translator demonstrates a complete MLOps lifecycle applied to a challenging
low-resource language scenario under real hardware constraints. Every rubric component
is implemented: Airflow orchestration, DVC data versioning, MLflow experiment tracking,
Prometheus/Grafana monitoring, FastAPI serving, Docker Compose packaging, and GitHub
Actions CI. The system deliberately prioritises \textbf{MLOps best practices over raw
model accuracy} --- the architecture is designed to detect performance degradation,
trigger automated retraining, and promote improved models to production without manual
intervention.

The no-cloud constraint, rather than being a limitation, motivated engineering decisions
(LoRA over full fine-tuning, Whisper-tiny over Whisper-small, dictionary fallback) that
are consistent with production MLOps for resource-constrained environments and are
directly supported by recent research literature.

% ──────────────────────────────────────────────────────────────────────────────
\begin{thebibliography}{9}

\bibitem{springer2024}
  Exploration of Whisper fine-tuning strategies for low-resource ASR.
  \textit{Journal on Audio, Speech, and Music Processing}, Springer, 2024.
  \url{https://link.springer.com/article/10.1186/s13636-024-00349-3}

\bibitem{lora2021}
  Hu, E. et al.\ LoRA: Low-Rank Adaptation of Large Language Models.
  \textit{arXiv:2106.09685}, 2021.

\bibitem{mdpi2025}
  LoRA-INT8 Whisper: A Low-Cost Framework for Edge Devices.
  \textit{MDPI Sensors}, 2025.
  \url{https://www.mdpi.com/1424-8220/25/17/5404}

\bibitem{peft}
  HuggingFace PEFT Documentation.
  \url{https://huggingface.co/docs/peft}

\bibitem{reykunyu}
  Reykunyu Na'vi Dictionary API.
  \url{https://reykunyu.lu}

\end{thebibliography}

\end{document}
